\documentclass[a4paper,11pt]{article}

\usepackage{amsmath,amssymb,bm}
\usepackage{physics}
\usepackage{geometry}
\usepackage{hyperref}
\usepackage{listings}
\usepackage{xcolor}
\geometry{margin=20mm}

\lstset{
  basicstyle=\ttfamily\small,
  keywordstyle=\color{blue},
  commentstyle=\color{gray},
  stringstyle=\color{teal},
  showstringspaces=false,
  frame=single,
  breaklines=true
}

\title{Fig4 (run\_fig4\_absD2\_vs\_absK2.py) のゲージ定義と\\
ゲージ不整合を絶対に再発させないための詳細メモ（再現可能版）}
\author{}
\date{\today}

\begin{document}
\maketitle

\section*{対象・前提}
\begin{itemize}
  \item 対象ファイルは \textbf{run\_fig4\_absD2\_vs\_absK2.py だけ}（あなたの要望どおり）。
  \item このメモは「\textbf{いつでも間違えたらこれを見れば直せる}」ことを目的に、\textbf{ゲージの定義}を数式で固定し、\textbf{コードとの対応}を（行番号付きで）明示する。
  \item 本メモは、あなたが採用している ``proxy（Airy フィルタ + 近似重み）'' の枠内で、\textbf{位相参照面（ゲージ）を厳密に整合させる}話に限定する
  （dyadic GF の係数実装の前段階の整合条件）。
\end{itemize}

\section*{結論（これだけ覚えれば良い）}
このスクリプトは
\begin{enumerate}
  \item 位相（光学長）に \(\{L_{\mathrm{cav,opt}}, z_{\mathrm{opt}}\}\) を使う（\textbf{prime/cavity-gauge}）
  \item 反射係数 \(r_b,r_t\) を \textbf{EML 界面基準（Gauge-B）}で計算する関数から取ってしまいがち
\end{enumerate}
なので、放置すると
\[
\boxed{
D = 1 - \underbrace{r_b^{(B)} r_t^{(B)}}_{\text{EML基準}}
\exp\!\left(i\frac{4\pi}{\lambda} \underbrace{L_{\mathrm{cav,opt}}}_{\text{prime基準}}\right)
}
\]
という \textbf{ゲージ混在}が発生する。

\bigskip
\noindent
これを再発させない唯一の方法は次のどちらか：
\begin{itemize}
  \item \textbf{(方法B)} すべてを Gauge-B にそろえて \(L_{\mathrm{eml,opt}}\) を指数に使う
  \item \textbf{(方法P)} 反射係数だけ位相補正して \(\{L_{\mathrm{cav,opt}}, r_b', r_t'\}\) にそろえる（``ゲージ分を取り直す''）
\end{itemize}
この文書と付属パッチは \textbf{(方法P)} を run\_fig4 の中だけで完結させている。

\section{ゲージ定義を数式で固定する}

\subsection{波長・面内パラメータ}
\begin{itemize}
\item 波長：\(\lambda\)（nm表記を想定）
\item 面内パラメータ：\(u\)（スクリプトは \(u\in[0,2]\) などの格子）
\item スクリプトの屈折率は、位相用に \textbf{実数}を使う：
\[
n_j^{(\mathrm{re})} \equiv \Re\{n_j\}
\]
\item スクリプトの光学座標は
\[
\Re\left(\frac{k_{z,j}}{k_0}\right)=\Re\sqrt{(n_j^{(\mathrm{re})})^2-u^2}
\]
で作る（吸収を位相座標から外す）。これは proxy の設計仕様。
\end{itemize}

\subsection{ゲージB（EML基準）}
\textbf{Gauge-B} は「EML を基準にしたゲージ」で、
\[
\{L_{\mathrm{eml,opt}}(u),\ r_b^{(B)}(\lambda,u),\ r_t^{(B)}(\lambda,u)\}
\]
を指す。

\begin{itemize}
\item \(L_{\mathrm{eml,opt}}(u)\)：
\[
\boxed{
L_{\mathrm{eml,opt}}(u)=
\Re\sqrt{(n_{\mathrm{EML}}^{(\mathrm{re})})^2-u^2}\,d_{\mathrm{EML}}
}
\]
\item \(r_b^{(B)}\)：\textbf{EML 下端界面（EML/EBL）側}を参照面とする、下側スタックへの見込み反射（複素）
\item \(r_t^{(B)}\)：\textbf{EML 上端界面（EML/ETL）側}を参照面とする、上側スタックへの見込み反射（複素）
\end{itemize}

Gauge-B で Fabry--Pérot 型分母を作るなら：
\[
\boxed{
D_B(\lambda,u)=
1-r_b^{(B)}(\lambda,u)\,r_t^{(B)}(\lambda,u)\,
\exp\!\left(i\frac{4\pi}{\lambda}L_{\mathrm{eml,opt}}(u)\right)
}
\]

\subsection{prime（cavity-gauge, プライム）}
\textbf{prime} は「ITO側の原点から ETL/cathode1 面まで」を光学座標で測るゲージで、
\[
\{L_{\mathrm{cav,opt}}(u),\ z_{\mathrm{opt}}(u),\ r_b'(\lambda,u),\ r_t'(\lambda,u)\}
\]
を指す。

\begin{itemize}
\item \(L_{\mathrm{cav,opt}}(u)\)：ITO起点から ETL/cathode1 面までの単一パス光学長（optcoord が返す）
\item \(z_{\mathrm{opt}}(u)\)：同じ原点から EML center までの光学座標
\end{itemize}

prime で分母を作るなら：
\[
\boxed{
D_P(\lambda,u)=
1-r_b'(\lambda,u)\,r_t'(\lambda,u)\,
\exp\!\left(i\frac{4\pi}{\lambda}L_{\mathrm{cav,opt}}(u)\right)
}
\]

\subsection{2つのゲージの関係（これが最重要）}
prime の全長は
\[
L_{\mathrm{cav,opt}}(u)=S_b(u)+L_{\mathrm{eml,opt}}(u)+S_t(u)
\]
と分解できる。ここで
\begin{equation}
S_b(u) =
\sum_{j\in\{\mathrm{ITO,pHTL,Rprime,HTL,EBL}\}}
\Re\sqrt{(n_j^{(\mathrm{re})})^2-u^2}\,d_j
\end{equation}

\begin{equation}
S_t(u) =
\Re\sqrt{(n_{\mathrm{ETL}}^{(\mathrm{re})})^2-u^2}\,d_{\mathrm{ETL}}
\end{equation}
は、prime の参照面と EML 界面参照面の ``ずれ（ゲージ差）'' の光学距離である。

\section{run\_fig4 の式（コードと同一の式で固定）}

\subsection{偏光ごとのフィルタ}
スクリプトは TE/TM（\(\alpha\in\{\mathrm{TE,TM}\}\)）ごとに
\[
F_\alpha(\lambda,u)=\frac{G_\alpha(\lambda,u)}{|D_\alpha(\lambda,u)|^2}
\quad\text{または}\quad
F_\alpha^{(c)}(\lambda,u)=\frac{G_\alpha(\lambda,u)}{D_\alpha(\lambda,u)}
\]
を作る（\(|D|^2\) 版と complex \(D\) 版）。

\subsubsection{分子 \(G\)}
コードの形（run\_fig4 では \texttt{phi\_sign=-1} がデフォルト）：
\[
\boxed{
G_\alpha(\lambda,u)=
1+|r_b^{(\alpha)}|^2
+2|r_b^{(\alpha)}|\cos\!\left(\arg(r_b^{(\alpha)})+\frac{4\pi}{\lambda}z_{\mathrm{opt}}(u)\right)
}
\]
これは次と等価：
\[
G_\alpha(\lambda,u)=
\left|1+r_b^{(\alpha)}(\lambda,u)\exp\!\left(i\frac{4\pi}{\lambda}z_{\mathrm{opt}}(u)\right)\right|^2
\]

\subsubsection{分母 \(D\)}
コードは振幅と位相に分解して計算しているが、厳密には
\[
\boxed{
D_\alpha(\lambda,u)=
1-r_b^{(\alpha)}(\lambda,u)\,r_t^{(\alpha)}(\lambda,u)\,
\exp\!\left(i\frac{4\pi}{\lambda}L_{\mathrm{cav,opt}}(u)\right)
}
\]
と同一である（\(\phi_{\mathrm{sign}}=-1\) のため）。

\bigskip
\noindent
\textbf{ここでの致命的な落とし穴}：
\begin{itemize}
\item \(L_{\mathrm{cav,opt}}\) を使うなら、\(r_b,r_t\) は prime 参照面の \(r_b',r_t'\) でなければならない。
\item もし \(r_b,r_t\) が Gauge-B 参照面の \(r_b^{(B)},r_t^{(B)}\) だと、指数は \(L_{\mathrm{eml,opt}}\) を使うべきで、混在する。
\end{itemize}

\section{混在を直す「ゲージ分を取り直す」処方箋（方法P）}

\subsection{狙い}
スクリプトは \(L_{\mathrm{cav,opt}}\) を使い続けたい（Fig4 表示ルーチンを変えたくない）ので、
\[
r_b^{(B)}, r_t^{(B)} \ \longrightarrow\ r_b', r_t'
\]
へ \textbf{参照面移送}を行い、
\[
r_b' r_t' e^{i\frac{4\pi}{\lambda}L_{\mathrm{cav,opt}}}
=
r_b^{(B)} r_t^{(B)} e^{i\frac{4\pi}{\lambda}L_{\mathrm{eml,opt}}}
\]
を恒等的に満たすようにする。

\subsection{位相補正（参照面移送）の定義}
run\_fig4 の proxy 位相定義（実数 \(n\)）に合わせて、
\begin{equation}    
r_b'(\lambda,u) = r_b^{(B)}(\lambda,u)\exp\!\left(-i\frac{4\pi}{\lambda}S_b(u)\right)
\end{equation}
\begin{equation}
r_t'(\lambda,u) = r_t^{(B)}(\lambda,u)\exp\!\left(-i\frac{4\pi}{\lambda}S_t(u)\right)
\end{equation}

とする。

このとき
\[
r_b' r_t' e^{i\frac{4\pi}{\lambda}L_{\mathrm{cav,opt}}}
=
r_b^{(B)} r_t^{(B)}
e^{i\frac{4\pi}{\lambda}(L_{\mathrm{cav,opt}}-S_b-S_t)}
=
r_b^{(B)} r_t^{(B)} e^{i\frac{4\pi}{\lambda}L_{\mathrm{eml,opt}}}
\]
なので、prime で計算しても Gauge-B で計算しても分母は一致する。

\section{コード対応（行番号付きで完全固定）}

本節は、\textbf{どこを見れば間違いを即修正できるか}のためのチェックリスト。

\subsection{(1) 位相ゲージ（prime）を作っている箇所}
\begin{itemize}
\item \texttt{zopt\_Lcavopt} / \texttt{zopt\_Lcavopt\_u} の import：
\begin{lstlisting}
(L32) from cavity_common.optcoord import zopt_Lcavopt
(L34) from cavity_common.optcoord_angle import zopt_Lcavopt_u
\end{lstlisting}

\item 実際に呼んで \(z,L\) を得る箇所：
\begin{lstlisting}
(L138) z0, L0 = zopt_Lcavopt(nmap, d)           # u=0
(L144) z_u, L_u = zopt_Lcavopt_u(nmap, d, u)    # u!=0
\end{lstlisting}

\textbf{ここで得る \(L\) は prime の \(L_{\mathrm{cav,opt}}\)。}
\end{itemize}

\subsection{(2) 分母 \(D\) の定義（prime を仮定していることの確定）}
\begin{itemize}
\item \(|D|^2\) 版の位相：
\begin{lstlisting}
(L102) dPhi = (-phi_sign*phib) + (-phi_sign*phit) + (4*pi*Lcavopt/lambda_nm)
\end{lstlisting}
\item complex \(D\)：
\begin{lstlisting}
(L113) D = 1.0 - sqrtRbRt*np.exp(1j*dPhi)
\end{lstlisting}

\textbf{この 2 行が「指数に \(L_{\mathrm{cav,opt}}\) を入れている」確定点。}
\end{itemize}

\subsection{(3) 反射係数が Gauge-B で来る（混在の入り口）}
このスクリプトが反射係数を取る箇所：
\begin{lstlisting}
(L135) rb0, rt0 = rb_rt_cavity_side(nmap, d, lam)             # u=0
(L145) rb_te, rt_te = rb_rt_cavity_side_u(nmap, d, lam, u, "TE")
(L146) rb_tm, rt_tm = rb_rt_cavity_side_u(nmap, d, lam, u, "TM")
\end{lstlisting}

\noindent
これらは（設計上）\textbf{EML界面参照（Gauge-B）}の \(r_b^{(B)},r_t^{(B)}\) と解釈すべきで、
prime 分母に直に入れると混在になる。

\subsection{(4) 混在を殺す処置：\_gauge\_retake を必ず通す}
本パッチが追加した ``ゲージ分の取り直し'' の本体：

\begin{lstlisting}
(L079-L094) def _gauge_retake_rb_rt_prime(...):
(L082-L088) S_b = sum_{ITO,pHTL,Rprime,HTL,EBL} Re(kz/k0)*d
(L090)       S_t = Re(kz_ETL/k0)*d_ETL
(L092)       phase_b = exp(-i*4*pi*S_b/lambda)
(L093)       phase_t = exp(-i*4*pi*S_t/lambda)
(L094)       return rb*phase_b, rt*phase_t
\end{lstlisting}

そして \textbf{反射係数を取った直後}に必ず通す：

\begin{lstlisting}
(L135-L140) u=0 branch:
  rb0, rt0 = rb_rt_cavity_side(...)
  rb0, rt0 = _gauge_retake_rb_rt_prime(lam, 0.0, rb0, rt0)

(L145-L153) u!=0 branch:
  rb_te, rt_te = rb_rt_cavity_side_u(...)
  rb_tm, rt_tm = rb_rt_cavity_side_u(...)
  rb_te, rt_te = _gauge_retake_rb_rt_prime(lam, u, rb_te, rt_te)
  rb_tm, rt_tm = _gauge_retake_rb_rt_prime(lam, u, rb_tm, rt_tm)
\end{lstlisting}

\paragraph{再発防止の鉄則}
\begin{itemize}
  \item \textbf{D の指数に \(L_{\mathrm{cav,opt}}\) を使うなら、必ず retake 済みの \(r_b',r_t'\) を使う。}
  \item 逆に retake を外したいなら、\textbf{指数を \(L_{\mathrm{eml,opt}}\) に変えて Gauge-B に統一}する。
\end{itemize}

\section{「間違えたかも？」を一発で検出する数値チェック}
\subsection{チェックの目的}
本当に
\[
r_b' r_t' e^{i4\pi L_{\mathrm{cav,opt}}/\lambda}
=
r_b^{(B)} r_t^{(B)} e^{i4\pi L_{\mathrm{eml,opt}}/\lambda}
\]
が（数値的に）一致しているかを \textbf{毎回自動で確認}する。

\subsection{チェック手順（run\_fig4 にそのまま埋め込み可能）}
例えば \(u\neq 0\) の枝で、retake 前後の量を保持し、以下を評価する：

\begin{lstlisting}
# before retake
rbB_te, rtB_te = rb_rt_cavity_side_u(...)

# after retake
rbP_te, rtP_te = _gauge_retake_rb_rt_prime(lam, u, rbB_te, rtB_te)

# build Lemlopt(u) in the same proxy convention
Lemlopt = _kz_re(float(np.real(nmap["EML"])), u) * float(d["EML"])

lhs = rbP_te * rtP_te * np.exp(1j * (4*np.pi*L_u/lam))       # prime form
rhs = rbB_te * rtB_te * np.exp(1j * (4*np.pi*Lemlopt/lam))   # gauge-B form
err = np.max(np.abs(lhs - rhs))
print("Gauge consistency err =", err)
\end{lstlisting}

\noindent
\textbf{期待値}：double 精度なら \(\sim 10^{-12}\) 程度（計算順序で多少上下）。

\section{よくある事故パターン（ここを見れば必ず直る）}

\subsection*{事故1：\(S_b\) に含める層を間違える}
\textbf{原因}：prime 原点から EML 下端までの積分に入る層集合がずれる。\\
\textbf{対策}：run\_fig4 の \texttt{\_gauge\_retake\_rb\_rt\_prime}（L82--L88）にある層集合
\[
\{\mathrm{ITO,pHTL,Rprime,HTL,EBL}\}
\]
が、あなたの stack 定義（optcoord が使う順序）と一致していることを毎回確認する。

\subsection*{事故2：\(S_t\) を ETL 以外も含めてしまう}
\textbf{原因}：prime の終点が ETL/cathode1 面であるのに、cathode1 の位相まで入れてしまう。\\
\textbf{対策}：このスクリプトの prime 終点は \textbf{ETL終端}で固定。よって
\[
S_t=\text{ETL のみ}
\]
（L90）が正しい。

\subsection*{事故3：位相用の \(k_z\) を複素屈折率で作ってしまう}
\textbf{原因}：proxy の位相座標が「実数 \(n\)」で設計されているのに、
複素 \(n\) を突っ込むと \(S_b,S_t\) が別物になり一致性が崩れる。\\
\textbf{対策}：このスクリプトは \textbf{実数 \(n\)}で \(\Re\sqrt{n^2-u^2}\) を使う。
それを \texttt{\_kz\_re}（L72--L76）で固定している。

\subsection*{事故4：retake を TE にだけ適用、TM を忘れる}
\textbf{対策}：L148--L149 で \textbf{TE/TM それぞれ}に適用していることを確認。

\section*{まとめ（最短の再発防止ルール）}
\begin{enumerate}
  \item \textbf{分母の指数}に何を使っているかでゲージが決まる。\\
  \(\exp(i4\pi L_{\mathrm{cav,opt}}/\lambda)\Rightarrow\) prime なので \(r_b',r_t'\)。\\
  \(\exp(i4\pi L_{\mathrm{eml,opt}}/\lambda)\Rightarrow\) Gauge-B なので \(r_b^{(B)},r_t^{(B)}\)。
  \item run\_fig4 は \(L_{\mathrm{cav,opt}}\) を使う。よって \textbf{retake が必須}。
  \item retake は
  \[
  r_b' = r_b^{(B)}e^{-i4\pi S_b/\lambda},\quad
  r_t' = r_t^{(B)}e^{-i4\pi S_t/\lambda}
  \]
  で、\(S_b\) の層集合と \(S_t=\)ETL のみを絶対に間違えない。
  \item 間違えたかも？と思ったら、\textbf{lhs--rhs の一致チェック（err）}を出す。
\end{enumerate}

\end{document}
